% Part: first-order-logic
% Chapter: natural-deduction
% Section: soundness-identity

\documentclass[../../../include/open-logic-section]{subfiles}

\begin{document}

\olfileid[pt]{fol}{ntd}{sid}

\olsection{Correção com \usetoken{S}{identity}}

\begin{prop}
A dedução natural com regras para $\eq$ é correta.
\end{prop}

\begin{proof}
Qualquer !!{formula} da forma $\eq[t][t]$ é válida, pois
para toda !!{structure}~$\Struct M$, $\Sat{M}{\eq[t][t]}$. (Note que
assumimos que o termo $t$ é fechado, isto é, não contém variáveis,
de modo que as atribuições de variáveis são irrelevantes).

Suponha que a última inferência em uma !!{derivation} é \Elim{\eq},
isto é, a !!{derivation} tem a seguinte forma:
\begin{prooftree}
  \AxiomC{$\Gamma_1$}
  \RightLabel{$\delta_1$}
  \DeduceC{$\eq[t_1][t_2]$}
  \AxiomC{$\Gamma_2$}
  \RightLabel{$\delta_2$}
  \DeduceC{$!A(t_1)$}
  \RightLabel{\Elim{\eq}}
  \BinaryInfC{$!A(t_2)$}
\end{prooftree}
As premissas $\eq[t_1][t_2]$ e $!A(t_1)$ são !!{derive}das a partir de
hipóteses !!{undischarged}s~$\Gamma_1$ e $\Gamma_2$, respectivamente.
Queremos mostrar que $!A(t_2)$ segue de $\Gamma_1 \cup \Gamma_2$.
Considere uma !!{structure}~$\Struct{M}$ com $\Sat{M}{\Gamma_1 \cup
  \Gamma_2}$. Pela hipótese de indução, $\Sat{M}{!A(t_1)}$ e
$\Sat{M}{\eq[t_1][t_2]}$. Portanto, $\Value{t_1}{M} = \Value{t_2}{M}$. Seja
$s$ qualquer atribuição de variáveis e $m = \Value{t_1}{M} = \Value{t_2}{M}$. Por
\olref[fol][syn][ext]{prop:ext-formulas}, $\Sat{M}{!A(t_1)}[s]$ se, e somente se,
$\Sat{M}{!A(x)}[\Subst{s}{m}{x}]$ se, e somente se, $\Sat{M}{!A(t_2)}[s]$. Como
$\Sat{M}{!A(t_1)}$, temos $\Sat{M}{!A(t_2)}$.
\end{proof}

\end{document}
